\section{An actual scalar primitive for the three-pivot extending error}
\label{sec:r12-three-pivot-error-primitive}

\paragraph{Status and attribution.}
\textbf{PROVED THIS ROUND, submitted for exact-source independent
review.}  The construction contributors are conservatively
\texttt{/root} and the R12 matrix agent
\texttt{/root/r11\_challenge}.  The exact G3F face formulas and
the frozen low-degree Steinberg identity in USR are inputs.
The result removes the unpunctured coefficient error of the
specified three-face defect.  It does not yet provide a primitive
of the remaining closed actual triple.

\subsection{The original analytic triangle behind the three faces}

Work on the entire ordinary three-pivot overlap $N_{adg}$ of G3F,
with arbitrary ambient rank $r\ge3$, $d_0=r-2$, and the original
closed metric curvature $R=c_1(\overline L)$.  To distinguish a
matrix minor from this curvature, put
\[
 P=ae-bd,\qquad Q=ah-bg,\qquad R_m=dh-eg.
 \tag{G3S1}
\]
The first two columns of an invertible matrix have rank two at
every analytic point.  Thus these three minors never vanish
simultaneously.  They satisfy the original analytic identity
\[
             aR_m=dQ-gP.
 \tag{G3S2}
\]
Define analytic functions and an analytic unit on $N_{adg}$ by
\[
 \alpha=P/a,\qquad \beta=Q/g,\qquad \gamma=R_m/g,
 \qquad \lambda=d/a.
 \tag{G3S3}
\]
Consequently
\[
        \alpha+\gamma=\lambda\beta,
        \qquad (\alpha,\beta,\gamma)\text{ has no common zero}.
 \tag{G3S4}
\]
This is a genuine analytic relation in the original field,
including its cancellation branch, not a freely imposed tropical
three-ray relation.  None of $\alpha,\beta,\gamma$ is required
to be a unit on all of $N_{adg}$.

On their common unit locus use the notation
\[
 x=v(\alpha),\quad y=v(\beta),\quad z=v(\gamma),\quad
 T=v(\lambda),\quad U=v(g/a).
 \tag{G3S5}
\]
The functions $T,U$ are actual unit logarithms on the whole
$N_{adg}$.  For unit logarithms $p,q$, use the original pair order
and write
\[
 e_{p,q}=(p\,d''q,q\,d'p),\qquad
 M_f=(0,0,f),\qquad
 \mathfrak F(p,q)=(1-\chi(q-p))e_{q,p-q},
 \qquad \mathfrak K(p,q)=\mathfrak F(p,q)-\mathfrak F(q,p).
 \tag{G3S6}
\]
The function $\chi$ is the same symmetric, constant-tail function
specified in G3F and SGI; its offset is denoted $\varepsilon>0$.
No separate extension of a singular term in (G3S6) is assumed.

\subsection{Reduction of the whole coefficient error}

The first-column valuation pairs of the Schur matrices in G3F are
exactly as follows, with signs ignored only because $|-1|=1$:
\[
\begin{array}{c|c}
\text{matrix}&\text{valuations of its two first-column entries}\\\hline
S_a&(x,y+U)\\
S_d&(x-T,z+U-T)\\
S_g&(z,y)\\
S_1=D_3((g/a)^{-1})S_a&(x,y)\\
J_{23}^{-1}S_1&(y,x)\\
J_{23}^{-1}S_g&(y,z)\\
S_1^{(2)}\text{ in the }d\text{-to-}g\text{ algorithm}&(x-T,z)\\
J_{23}^{-1}S_1^{(2)}&(z,x-T)\\
J_{23}^{-1}(J_{23}^{-1}S_g)&(z,y).
\end{array}
 \tag{G3S7}
\]
The last three rows are obtained by applying G3F22 to
$J_{12}^{-1}M$, whose first-column entries are $d,-a,g$.
This retains the alternate-chart identification in G3F27; it
contributes no extending error.

Let $E_{adg}=e_{ad}+e_{dg}-e_{ag}$ be the \emph{entire}
unpunctured actual pair defined by G3F29.  Substitution of the
full G3F17 and G3F23 errors, followed only by literal bilinearity,
gives on the common unit locus
\[
 E_{adg}=\mathcal E\,R^{d_0},\qquad
 \mathcal E=-e_{T,x}+\mathfrak K(x-T,z)
                         -\mathfrak K(x,y)+\mathfrak K(z,y).
 \tag{G3S8}
\]
This is not division or cancellation of $R^{d_0}$: the coefficient
pair on the right is computed from the displayed coefficient
formulas before the common closed curvature is wedged.

Here is a useful check on all outer-unit and GMU-gauge terms.
Write $A=v(a)$, $D=v(d)$, $G=v(g)$ and $p=v(P)$.
The three outside CS errors have sum $e_{D-G,A-D}$.
The two row-two GMU gauge terms have sum
$e_{G-D,p-D}-e_{G-A,p-A}$.  Since $T=D-A$, $U=G-A$ and
$x=p-A$, their total is
\[
 e_{U-T,T}+e_{U-T,x-T}-e_{U,x}=-e_{T,x}.
 \tag{G3S9}
\]
The remaining six $\mathfrak F$ terms are exactly the three
$\mathfrak K$ terms in (G3S8).  Thus neither a row-unit gauge
nor an outside CS correction has been dropped.

\subsection{The scalar formula on the whole unit locus}

Define the ordinary smooth function
\[
 H_\chi(t)=\int_0^t s\chi(s)\,ds.
 \tag{G3S10}
\]
It satisfies $H_\chi'(t)=t\chi(t)$ and, by symmetry of $\chi$,
\[
 H_\chi(t)+H_\chi(-t)=t^2/2.
 \tag{G3S11}
\]
There is a constant $h_-$ with $H_\chi(t)=h_-$ for
$t\le-\varepsilon$ and $H_\chi(t)=t^2/2-h_-$ for
$t\ge\varepsilon$.
An exact first-form product computation gives
\[
 \mathfrak K(p,q)=e_{q,p-q}+\partial_2M_{H_\chi(q-p)}.
 \tag{G3S12}
\]
Indeed $e_{q,p-q}+e_{p,q-p}=\partial_2M_{-(p-q)^2/2}$;
insert this in (G3S6) and use the chain rule in (G3S10).
This is a first-form identity with all two-colour derivatives,
not a claim that the weighted summands are individually closed.

On the common unit locus set
\[
 a_0=x-T-y,\qquad b_0=z-T-y,
 \qquad c_0=\min(a_0,0)=\min(b_0,0).
 \tag{G3S13}
\]
The equality of the two minima follows from the two analytic units
$\alpha/(\lambda\beta)$ and $\gamma/(\lambda\beta)$, whose sum
is one by (G3S4).  The entire actual Steinberg identity USR15 gives
$e_{b_0,a_0}=\partial_2M_{c_0^2/2}$ on all its valuation walls.
Inserting (G3S12) in (G3S8), the remaining unit-symbol part is
\[
 \begin{split}
 &-e_{T,x}+e_{z,x-T-z}-e_{y,x-y}+e_{y,z-y}\\
 &\hspace{10mm}=e_{b_0,a_0}-e_{b_0,b_0}
                                    -\partial_2M_{Tz}.
 \end{split}
 \tag{G3S14}
\]
All cross terms are retained: the last correction comes from
$e_{b_0,T}+e_{T,b_0}+e_{T+y,T}+e_{T,T+y}$.
Consequently the explicit scalar
\[
 F=\frac{c_0^2-b_0^2}{2}-Tz
       +H_\chi(z-x+T)-H_\chi(y-x)+H_\chi(y-z)
 \tag{G3S15}
\]
satisfies $\mathcal E=\partial_2M_F$ on the entire common unit
locus, including the complete nonproduct triangle and all its
intercept walls.  In particular no verification only on strict
Gauss chambers is being used.

\subsection{Actual extension through every zero minor}

The formula has a direct finite ordinary extension, not an assigned
value to an infinite valuation.  Put
\[
 B=1+\varepsilon+2(|T|+|U|)
 \tag{G3S16}
\]
and consider the ordinary opens
\[
 \begin{split}
 O_\alpha&=\{|\alpha|<e^{-B}|\gamma|\},\\
 O_\beta&=\{|\lambda\beta|<e^{-B}|\alpha|\},\\
 O_\gamma&=\{|\gamma|<e^{-B}|\alpha|\}.
 \end{split}
 \tag{G3S17}
\]
Their defining norms are continuous at every analytic point;
$B$ has a finite actual unit-valuation presentation.  These opens
contain the respective zero loci because of (G3S4).  The other
two entries are units on each such open.  Together with the common
unit locus they cover $N_{adg}$.

On $O_\alpha$, the ultrametric relation gives $z=T+y$ on the whole
ordinary open.  Where $\alpha\ne0$, both $z-x+T$ and $y-x$ are
below $-\varepsilon$, $b_0=c_0=0$, and (G3S15) simplifies to
\[
                 F_\alpha=-Tz+H_\chi(-T).
 \tag{G3S18}
\]
The right side is actual on all of $O_\alpha$.
On $O_\beta$, one has $x=z$ and $a_0=b_0=c_0$ on its unit part;
the two large-$y$ terms cancel as whole functions.  Thus
\[
                 F_\beta=-Tz+H_\chi(T)
                           \quad\hbox{on }O_\beta.
 \tag{G3S19}
\]
On $O_\gamma$, one has $x=T+y$, $a_0=c_0=0$ and
$b_0=z-x>\varepsilon+2|T|$.  Use both exact tails in (G3S11):
\[
 \begin{split}
 F&=-b_0^2/2-Tz+(b_0+T)^2/2-H_\chi(-T)\\
  &=-Tx+H_\chi(T).
 \end{split}
 \tag{G3S20}
\]
The constants $h_-$ cancel, so this is an exact scalar identity,
not a statement modulo a constant.  Define $F_\gamma$ by the last
expression on all of $O_\gamma$.

These explicit finite functions agree with (G3S15) on their whole
ordinary intersections with its domain.  Hence they glue to one
actual $F\in B^{0,0}(N_{adg})$.  No simultaneous-zero point or
pole has been silently added to the domain.

It remains important to check the pair identity at a zero minor
itself, rather than infer it by restriction injectivity.  The
positive $|U|$ buffer in (G3S16) puts every singular first-pivot
term in the original table (G3S7) in the corresponding constant
tail of its GB partition.  Inserting those constant tails into the
\emph{combined} GMU12--GMU15 and TML4 ordinary extensions gives
the following direct coefficient computations.
On $O_\alpha$, the $x$ terms cancel before any extension, and
\[
 \begin{split}
 \mathcal E
   &=-e_{T,x}+e_{z,x-T-z}-e_{y,x-y}+\mathfrak K(z,y)\\
   &=-\partial_2M_{Tz}+\partial_2M_{H_\chi(-T)}
                       =\partial_2M_{F_\alpha}.
 \end{split}
 \tag{G3S21}
\]
The second line contains only the two finite unit logarithms
$T,z$ and is the prescribed whole extension at $\alpha=0$.
On $O_\beta$, $x=z$ makes the entire two singular terms
$-\mathfrak K(x,y)+\mathfrak K(z,y)$ cancel, including their
partition derivatives.  The remaining expression is
\[
 \mathcal E=-e_{T,x}+\mathfrak K(x-T,x)
                =\partial_2M_{-Tx+H_\chi(T)}.
 \tag{G3S22}
\]
On $O_\gamma$, the two tails give
$\mathfrak K(x-T,z)=-e_{y,z-y}$ and
$\mathfrak K(z,y)=e_{y,z-y}$, since $x-T=y$.
They cancel as whole expressions.  Thus
\[
 \begin{split}
 \mathcal E&=-e_{T,x}-\mathfrak K(x,y)\\
   &=-\partial_2M_{Tx}+\partial_2M_{T^2/2-H_\chi(-T)}
                      =\partial_2M_{F_\gamma}.
 \end{split}
 \tag{G3S23}
\]
In each case the disappearing logarithm is eliminated by a literal
unit-logarithm identity on the ordinary zero neighborhood before
the original whole-call extension is applied.  No singular scalar
or pair has been extended separately, and no unpunctured equality
is deduced merely from a punctured equality.

\begin{theorem}[Whole-neighborhood coefficient-error primitive]
The scalar $F$ in (G3S15), (G3S18)--(G3S20) is actual on the entire
ordinary $N_{adg}$ and satisfies
\[
          E_{adg}=\partial_2(M_F R^{d_0})
                      \quad\hbox{on all of }N_{adg}.
 \tag{G3S24}
\]
Consequently the prescribed whole triple
\[
           \mathcal C_{adg}=T_{adg}+M_FR^{d_0}
 \tag{G3S25}
\]
is actual and $\partial_2$-closed on the punctured three-chart
overlap.  Equivalently, the actual extending triple
$F_{adg}=-M_FR^{d_0}$ solves the stronger unpunctured target
$\partial_2F_{adg}=-E_{adg}$ of G3F31.
\end{theorem}

\begin{proof}
Equations (G3S8)--(G3S15) prove (G3S24) on the common unit locus
using the complete actual Steinberg relation.  Equations
(G3S21)--(G3S23) verify exactly the same identity on all three
ordinary zero neighborhoods, in their prescribed whole extensions.
These opens cover.  Since $R$ is independently closed in both
colours, wedging the displayed scalar primitive with $R^{d_0}$
commutes with the actual boundary, without canceling that factor.
G3F30 then gives (G3S25).
\end{proof}

\paragraph{The remaining actual, not formal, task.}
An actual $H$ with $\mathcal C_{adg}=\partial_3H$, or more
generally an explicit relative third witness with a further actual
extending closed triple, has not been supplied by this theorem.
The closed triple in (G3S25) still contains every original native
word and cap term and EFC's product-torus integrals.  It is not
automatically an AGC cube or an arbitrary family to which PTH can
be applied.  A genuine product-presented tail certificate for that
\emph{whole} triple, or a direct complete higher identity, is still
required.  Therefore the general three-row block Green application
remains conditional on the remaining part of G3F31.

\paragraph{Exact inputs and boundaries.}
G3F: \path{c1574d77ea47992a02f5ebc80fce519dbad5927223f00622a1eddc33cbf8dd44}.
USR: \path{d7225979dee2fba900e590f725754ab4c59054caa674d1d2ad749b86b81b5acd}.
GMU: \path{a70b10906ebf4c3259df5d69131a3726f3b14aa8040f9df319a3e03bf8338fcb}.
TML: \path{327b233fe982c1692eba9603cfd50f0ff35932682fff182322343bc57ff6fddd}.
The actual triple and first-form conventions are their frozen CPT
conventions.  The low-degree Steinberg identity is a derived
interface to frozen GPR13, not claimed anew here.  The new result
is the full three-face error reduction and its all-zero-domain
scalar primitive.  It holds in every original residue
characteristic, for nonrational points and nonproduct base units,
with no field extension, new growth condition or probe-topology
substitution.  It does not close the full actual kernel or general
Strong G.
